% part: intuitionistic-logic
% chapter: propositions-as-types
% section: types

\documentclass[../../../include/open-logic-section]{subfiles}

\begin{document}

\olfileid{int}{pty}{typ}

\olsection{ټايپونه}

ثبوت ترم، مثلاً $(MN)$، پخپله په يکتا ډول هغه ثبوت
نۀ ټاکي چې دا ئې ثبوت ترم دے۔ ځکه آزاد متغيرونه، يعنې
د اکسيومونو ثبوت ترمونه، د اړوندو اکسيومونو !!{formula}s
نۀ ښيي۔ خو که دا په سياق~$\Gamma$ کښې وټاکو، نو
ثبوت \emph{په يکتا ډول} ټاکل کېږي۔

\begin{defn}
د ثبوت ترم~$N$ دپاره \emph{سياق} داسې مجموعه $\Gamma$
ده چې د~$N$ د هر آزاد متغير دپاره په $\Gamma$ کښې
د $x : !A$ کټ مټ يوه جوړه وي۔
\end{defn}

په ياد ولرئ چې تعريف دا شرط نۀ ږدي چې $\Gamma$
دې $x:!A$ \emph{يوازې} د~$N$ د آزادو متغيرونو دپاره
ولري۔ مثلاً $\Gamma = \{x:!A, y:!B, z:!C\}$
د $\pair{x}{y}$ دپاره يو سياق دے۔

\begin{defn}\ollabel{defn:type}
  فرض کړئ $\Gamma$ د~$N$ دپاره يو سياق دے۔ د~$N$
  \emph{ټايپ} د~$\Gamma$ په نسبت په استقرايي ډول داسې
  تعريفېږي:
  \begin{enumerate}
  \item د $x$ ټايپ $!A$ هغه وخت او يوازې هغه وخت دے
  چې $x:!A \in \Gamma$ وي۔
  \item که $N$ د $x:!A, \Gamma$ په نسبت ټايپ~$!B$
  ولري، نو د $\lambd[\typeof{x}{!A}][N]$ ټايپ $!A \lif !B$ دے۔
  \item که $N$ ټايپ $!A \lif !B$ او $M$ ټايپ~$!A$
  ولري، نو $(NM)$ ټايپ~$!B$ لري۔
  \item که $N$ ټايپ $!A$ او $M$ ټايپ~$!B$ ولري،
  نو $\pair{N}{M}$ ټايپ $!A \land !B$ لري۔
  \item که $N$ ټايپ~$!A_1 \land !A_2$ ولري، نو
  $\proj{i}{N}$ ټايپ~$!A_i$ لري۔
  \item که $N$ ټايپ $!A$ ولري، نو
  $\inj[!B]{i}{N}$ ټايپ $!A \lor !B$ لري که $i=1$
  وي، او ټايپ $!B \lor !A$ لري که $i=2$ وي۔
  \textbf{د سرچينې سمون (OLCMP-219، د210تکرار):}
  د جوړوونکي آرګومېنټ په اصل کښې د فرض شوي ترم پر
  ځای بل، بې‌تعريفه فورمولي نوم و؛ هماغه فرض شوي ترم
  ته سم شوے۔ شاخصونه او د دواړو ټايپونو ترتيب عين دي۔
  \item که $M$ ټايپ $!A \lor !B$ ولري، او $N_1$
  د $x:!A, \Gamma$ په نسبت ټايپ $!C$ او $N_2$
  د $y:!B, \Gamma$ په نسبت ټايپ~$!C$ ولري، نو
  $\dcase{M}{x}{N_1}{y}{N_2}$ ټايپ~$!C$ لري۔
  \textbf{د سرچينې سمون (OLCMP-220):} د حالتونو په
  ترم کښې تړونکي متغيرونه د دواړو مقدماتي سياقونو
  له عين متغيرونو سره برابر شوي؛ اصل بل شاخص لرونکے
  جوړه کاروله چې په هغو سياقونو کښې نۀ وه۔
  \item که $N$ ټايپ~$\lfalse$ ولري، نو
  $\abort{!A}{N}$ ټايپ~$!A$ لري۔
  \end{enumerate}
\end{defn}

\begin{prop}
هر ثبوت ترم~$N$ د~$N$ دپاره د سياق~$\Gamma$ په نسبت
تر ډېره يو ټايپ لري۔
\textbf{د سرچينې سمون (OLCMP-221):} اصل د هر خام
ترم د ټايپ وجود او يکتايي يو ځای ادعا کوي، خو د همدې
برخې ځان‌تطبيق بېلګه هېڅ ټايپ نۀ لري۔ صحيح دعوه دا ده
چې که د ترم ټايپ موجود وي، يکتا دے؛ يوازې د سياق
ټاکل د خام ترم د ټايپ وجود نۀ تضمينوي۔
\end{prop}

\begin{proof}
  پر~$N$ په استقرا۔
\end{proof}

که د $N$ ټايپ د $\Gamma$ په نسبت~$!A$ وي،
نو $\Gamma \Sequent N : !A$ ليکو۔ د \olref{defn:type}
مواد بايد اشنا ښکاره شي: دا کټ مټ د \Log{tN2i} قواعد
دي، يوازې په يوه لږ توپير سره؛ سياق~$\Gamma$ په ټوله
لړۍ کښې هماغه دے۔ دا تعريف د \Log{tN3i} په حساب کښې
د \olref{tab:tN2ip} د قواعدو په وسيله ليکلے شو۔

\subfile{rules-tN3}

ثبوت ترمونه د $\lambd$-حساب د يوه ډول ترمونه دي؛
يعنې \emph{له حاصل‌ضربونو، مجموعو او تش ټايپ سره
ټايپ‌لرونکے $\lambd$-حساب}۔ دوديز $\lambd$-حساب يوازې
له متغيرونو، \emph{$\lambd$-تجريدونو}~$\lambd[x][N]$
او تطبيق $(MN)$ څخه جوړ ترمونه لري، او د ټايپ‌لرونکي
$\lambd$-تجريد $\lambd[\typeof{x}{!A}][N]$ د~$!A$
نښه پکې نشته۔ $\lambd$-حساب د محاسبې داسې موډل دے
چې ټيورېنګ بشپړ دے؛ يعنې هره جزوي محاسبه کېدونکې
تابع پکې د يو ترم په وسيله تمثيلېدلے شي۔ د $\lambd$-حساب
هغه ډول چې ثبوت ترمونه او د ټايپ ټاکلو قواعد ئې ورکوي،
د محاسبې محدود موډل دے، خو بنسټيزې مفکورې هماغه دي۔
د ټايپ ټاکنه ئې محدودوي: يوازې \emph{صحيح ټايپ‌لرونکي}
ترمونه، يعنې هغه ثبوت ترمونه چې د سياق~$\Gamma$ په
نسبت ټايپونه لري، روا ترمونه دي۔ مثلاً $\lambd[x][(xx)]$
په دوديز بې‌ټايپه $\lambd$-حساب کښې روا ترم دے، خو
$\lambd[\typeof{x}{!A}][(xx)]$ د هېڅ $!A$ دپاره
صحيح ټايپ نۀ لري۔ ځکه د ټايپ قواعد دا شرط ږدي چې د
$(xx)$ د ټايپ لرلو دپاره لومړنے~$x$ بايد ټايپ $!A \lif !B$
او دويم $x$ ټايپ~$!A$ ولري؛ او دا ناشونې ده، ځکه
$!A$ د~$!A \lif !B$ يوه حقيقي فرعي !!{formula} ده۔

په وجداني ډول ټايپونه د هغو شيانو ډول ګڼلے شو چې
د هماغه ټايپ ترم ئې نوم اخلي۔ نو د ټايپ $!A \land !B$
ترم داسې جوړه ټاکي چې لومړۍ برخه ئې ټايپ~$!A$ او
دويمه برخه ئې ټايپ~$!B$ ولري؛ يعنې د حاصل‌ضرب~$!A \times !B$
!!a{element}۔ د ټايپ $!A \lif !B$ ترم له هغو شيانو
چې ټايپ $!A$ لري، هغو شيانو ته تابع ده چې ټايپ~$!B$
لري۔ د ټايپ $!A \lor !B$ ترم يا د ټايپ $!A$ يو شے
دے يا د ټايپ~$!B$ يو شے، او دا معلومات هم ورسره وي
چې کوم يو دے۔ دا د دوو مجموعو $!A$ او $!B$ د
بې‌اشتراکه اتحاد په شان دے، چې داسې تعريفېږي:
$\{\tuple{i,x} : i = 1 \text{ يا } 2, x \in !A \text{ که } i = 1, x \in !B \text{ که } i=2\}$۔
ټايپ $\lfalse$ تشه مجموعه ده؛ يعنې که يو ترم دا ټايپ
ولري، د هېڅ شي نوم نۀ شي اخيستلے۔ څرنګه چې طبيعي
اشتقاق، د \Log{tN3i} په ګډون، سازګار دے، د~$\lfalse$
!!{proof} هله شته چې خلاص فرضونه موجود وي؛ نو داسې
تړلے ثبوت ترم، يعنې بې له آزادو متغيرونو ترم، نشته
چې ټايپ~$\lfalse$ ولري۔

د ټايپ‌لرونکي $\lambd$-حساب عملګرونه داسې ترمونه
جوړوي چې په وجداني ډول د سم ډول شيانو نوم اخلي، په
دې شرط چې هغه ترمونه چې عملګر پرې پلي کېږي، د ټاکلو
شيانو نوم واخلي۔ مثلاً دا خبره چې ``که $N_1 : !A$ او
$N_2 : !B$ وي، نو $\pair{N_1}{N_2} : !A \land !B$''
دا وايي: که $N_1$ د ټايپ~$!A$ د يو شي او $N_2$
د ټايپ~$!B$ د يو شي نوم اخلي، نو $\pair{N_1}{N_2}$
د ټايپ $!A \land !B$ د يو شي نوم اخلي۔

\end{document}
